\section{The Collatz problem: an engine on bounded fuel}\label{sec:collatz}

Every mask so far has read a hard problem through a \emph{surrogate} engine. Collatz is the
exception: here the engine is a \emph{literal} object on $\N$. We translate the conjecture into the
language of the engine and the rope, locate exactly where it breaks the descent law, and drive the
green mechanics all the way to a named rope law --- which, in its universal form, we
machine-\emph{refute} in this very section. That refutation is not incidental: the law had been
accepted as the fourth boundary of the first cause \axiom{}, and its fall is the section's central
event.

\subsection{The map, and where it breaks the descent law}

\begin{definition}[accelerated Collatz map]\label{def:collatzMap}
The accelerated Collatz map $T\colon\N\to\N$ folds the forced halving after an odd step into itself:
\begin{equation}\label{eq:collatz-map}
T(n)=\begin{cases}n/2,& n\equiv 0\ (\mathrm{mod}\ 2),\\[2pt] (3n+1)/2,& n\equiv 1\ (\mathrm{mod}\ 2).\end{cases}
\end{equation}
The even step is a descent by exactly two (the edge case $A=2$, not strict); the odd step is a
\emph{fuel injection} --- the ``$+1$'' in $3n+1$ raises the height. \leantag{EuclidsPath.Collatz.Engine.T}
\end{definition}

Euclid's engine halts with a guarantee because every step is a strict descent, whence
\code{no\_infinite\_descent} forces $H(t)<H_0/A^{t}<1$. Collatz enjoys no such law.

\begin{theorem}[odd step ascends]\label{thm:odd_accel_ascends}
\green{}\ For odd $n\ge 1$ the accelerated step strictly raises the height: $n<\tfrac{3n+1}{2}=T(n)$.
\leantag{odd\_accel\_ascends}
\end{theorem}

\begin{theorem}[odd step is not a descent]\label{thm:not_descent_on_odd}
\green{}\ For odd $n\ge 1$ the step is not a descent, $\neg\,(T(n)<n)$. Hence the height along an orbit
is not strictly decreasing, and \code{no\_infinite\_descent} does not apply.
\leantag{not\_descent\_on\_odd}
\end{theorem}

\emph{Proof idea.} Both are one-line arithmetic on $\N$. Numerically the height wanders upward by
four to five orders of magnitude before descending (peak/start $\approx 5.4\times 10^{4}$ at
$n=159487$); there is no strict monotone Lyapunov function. This is the exact boundary between the
proven halting of Euclid's engine and the openness of Collatz: the fuel injection temporarily lifts
the height, so the strict form of the second law is broken.

\begin{theorem}[no monotone height]\label{thm:no_monotone_height}
\green{}\ Ascending positions are unbounded: for every $N$ there is $n>N$ with $n\equiv 1\ (\mathrm{mod}\ 4)$
and $n<(3n+1)/2$ (witness $n=4N+5$). Hence no monotone $\log_2$-height exists.
\leantag{no\_monotone\_height}
\end{theorem}

Averaged, the fuel is net-consumed --- the mean geometric drift per step is $\sqrt{3}/2\approx
0.864<1$, and halvings/triplings $=2.016$ exceeds the threshold $\log_2 3=1.585$. But this is a law
of the \emph{average}, while the conjecture demands that \emph{every} trajectory halts. Our
higher-energy incompatibility requires strict descent at each step; Collatz is an engine on bounded
fuel by energy, but not an EPMI engine.

\subsection{Tug of war: a multiplicative window budget}

We read the dynamics without logarithms. The engine advances by at most one binary rank per step;
the rope pulls back exactly one rank (halving) or two ($n\equiv 0\ \mathrm{mod}\ 4$). Descent is
then a bookkeeping inequality between the counts of even and odd steps in a window.

\begin{theorem}[engine advances at most one rank]\label{thm:engine_at_most_one_rank}
\green{}\ One move of the engine advances by at most a rank: $T(n)\le 2n$ for all $n$.
\leantag{engine\_at\_most\_one\_rank}
\end{theorem}

\begin{theorem}[window budget]\label{thm:window_budget}
\green{}\ Over a window of $k$ steps, of which $t=\mathrm{oddCount}(k,n)$ are odd and
$e=\mathrm{evenCount}(k,n)$ are even, the per-step bounds multiply under any order of steps:
\begin{equation}\label{eq:window-budget}
2^{e}\cdot T^{k}(n)\;\le\;2^{t}\cdot n.
\end{equation}
This is pure $\N$-arithmetic: the floor-divisions are absorbed by the precision of the per-step
bounds, and no logarithms are needed. \leantag{window\_budget}
\end{theorem}

\begin{theorem}[window descends]\label{thm:window_descends}
\green{}\ If the rope wins the window strictly ($t<e$) and $n\ge 1$, then $T^{k}(n)<n$.
\leantag{window\_descends}
\end{theorem}

\begin{proof}
From \cref{eq:window-budget}, $t<e$ gives $2^{t+1}\le 2^{e}$, hence $2\cdot T^{k}(n)\le n$, so
$T^{k}(n)<n$.
\end{proof}

The exact threshold $e>t$ --- not the average-implausible $e>2t$ --- is precisely the one that the
green $\times 2$ bound converts into a descent. Renouncing $\log_2 3$ costs a $\sim 1.6\%$ margin
(true descent needs only $e>0.585\,t$, but that is already logarithms, not core Lean).

\begin{definition}[rope counting law]\label{def:ropeCountingLaw}
For $n$, the rope dominance law asserts that above the absorbing cycle every position opens a window
with an even surplus:
\begin{equation}\label{eq:rope-law}
\mathrm{RopeCountingLaw}(n)\ :\equiv\ \forall j,\ \bigl(2<T^{j}(n)\Rightarrow \exists k,\ \mathrm{oddCount}(k,T^{j}(n))<\mathrm{evenCount}(k,T^{j}(n))\bigr).
\end{equation}
\leantag{RopeCountingLaw}
\end{definition}

\begin{theorem}[the counting law forces halting]\label{thm:reaches_one_of_countingLaw}
\green{}\ If $n\ge 1$ and $\mathrm{RopeCountingLaw}(n)$ holds, then $\exists K,\ T^{K}(n)=1$.
\leantag{reaches\_one\_of\_countingLaw}
\end{theorem}

\emph{Proof idea.} The ladder is entirely green: $\mathrm{RopeCountingLaw}(n)$ feeds the window
budget (\cref{thm:window_budget}), which yields a value-descent law, which by induction on the value
ceiling reaches $1$. This is a \emph{conditional} theorem about a specific $n$ --- ``give the law for
this $n$ and you get halting''. What is no longer permitted is to hand the premise to all $n$ at
once; the universal form is refuted below.

\subsection{The forged refutation: the climbing orbit of $27$}

The rope law demands an even surplus at \emph{every} position above the cycle, including the very
start $j=0$. The orbit of $27$ --- the most famous number in Collatz folklore --- climbs to a peak of
$4616$ before surrendering, and it never opens such a window from the start: while the number climbs,
the steps are odd-heavy, and $\mathrm{evenCount}-\mathrm{oddCount}$ in the prefix windows never turns
positive.

\begin{theorem}[the orbit of $27$ counted out]\label{thm:counts_at_70_at_27}
\green{}\ The orbit of $27$ reaches one in $70$ accelerated steps:
\begin{equation}\label{eq:counts-27}
\mathrm{oddCount}(70,27)=41,\qquad \mathrm{evenCount}(70,27)=29,\qquad T^{70}(27)=1.
\end{equation}
Forty-one moves of the engine against only twenty-nine pulls of the rope; the rope trails by twelve
overall and never once leads from the start. \leantag{counts\_at\_70\_at\_27}
\end{theorem}

\emph{Proof idea.} The kernel closes the prefix up to $k=70$ by direct computation: every prefix window
has $\mathrm{even}\le\mathrm{odd}$, and at $k=70$ exactly $41$ against $29$, finishing at $1$. The tail
falls into the cycle $1\to 2\to 1$, where steps go strictly in pairs; a cycle lemma, via the additivity
of the counters, freezes the deficit of $-12$ forever, so no window from the start ever flips it.

\begin{theorem}[the rope law fails at $27$]\label{thm:not_ropeCountingLaw_27}
\green{}\ The rope law is false for $n=27$: $\neg\,\mathrm{RopeCountingLaw}(27)$. From $j=0$ (value $27>2$)
no window gives a rope surplus. \leantag{not\_ropeCountingLaw\_27}
\end{theorem}

\begin{theorem}[the universal law is refuted]\label{thm:ropeLaw_universal_refuted}
\green{}\ The universal form of the law is false: $\neg\,\forall n\,(1\le n\Rightarrow\mathrm{RopeCountingLaw}(n))$,
with witness $n=27$ (\cref{thm:not_ropeCountingLaw_27}). A decree boundary on this law is impossible.
\leantag{ropeLaw\_universal\_refuted}
\end{theorem}

The beautiful metaphor ``the fuel runs out, the engine is dragged back to the start'' is true in
\emph{outcome} ($27$ does reach $1$) but not for the strong \emph{stepwise} reason the law asserted.
And $27$ is not alone: the harness finds a whole family of climbing violators ($27,31,41,47,54,55,\dots$).

\subsection{A non-descending orbit is a perpetual engine}

\begin{theorem}[a non-descending orbit is a perpetual engine]\label{thm:nonDescendingOrbit_is_perpetual_engine}
\green{}\ For $n\ge 3$ a non-descending orbit ($n\le T^{k}(n)$ for all $k$) loses no window to the rope
($\forall k,\ \mathrm{evenCount}(k,n)\le\mathrm{oddCount}(k,n)$) and never halts
($\forall K,\ T^{K}(n)\neq 1$). \leantag{nonDescendingOrbit\_is\_perpetual\_engine}
\end{theorem}

\emph{Proof idea.} Non-descent forces the window budget \cref{eq:window-budget} to keep the even
count from ever exceeding the odd count; an infinite feed with not a single lost window is precisely a
perpetual engine, and \cref{thm:window_descends} shows it cannot reach $1$.

\begin{theorem}[non-halting carries a perpetual engine]\label{thm:nonHalting_carries_perpetual_engine}
\green{}\ If an orbit never reaches $1$, then its tail from the global-minimum position is a non-descending,
never-halting orbit --- exhibited explicitly, not by contradiction.
\leantag{nonHalting\_carries\_perpetual\_engine}
\end{theorem}

The natural numbers are well-ordered, so every orbit has a global minimum
(\code{exists\_min\_position}, the standard triple); the tail from it never dips below its start.
Read aloud: \textbf{every Collatz counterexample carries a perpetual engine within it.} The
conjecture does not \emph{resemble} a ``no perpetual engine'' statement --- for the map $T$ it
\emph{is} one, word for word: an orbit reaches $1$ if and only if it has no perpetual tail.

\subsection{Verification, not derivation}

To \emph{learn the cause from inside} --- to derive the universal rope law by an internal proof ---
would mean extracting an infinite fact about the whole number line without leaving it. We formalise
the attempt as a self-grounding that crosses its own boundary.

\begin{definition}[internalised Collatz ground]\label{def:internalisedCollatzGround}
A structure with two fields: the \emph{ground}, the law itself $\forall n\ge 1,\ \mathrm{RopeCountingLaw}(n)$,
and a witness that it was derived by crossing the boundary,
$\neg\,\forall n\ge 1,\ \mathrm{RopeCountingLaw}(n)$. The abbreviation
\code{InternalKnowledgeOfCollatzCause} identifies internal knowledge of the cause with this
self-grounding. \leantag{InternalisedCollatzGround}
\end{definition}

\begin{theorem}[the Collatz cause is unknowable]\label{thm:collatzCause_unknowable}
\green{}\ Internal knowledge of the Collatz first cause is impossible (with no axioms at all):
$\neg\,\mathrm{InternalKnowledgeOfCollatzCause}$. \leantag{collatzCause\_unknowable}
\end{theorem}

\begin{proof}
The two fields of \cref{def:internalisedCollatzGround} contradict each other directly:
$\texttt{beyondOwnHorizon}\ \texttt{ground}:\mathrm{False}$.
\end{proof}

After the decree's fall this holds a fortiori and by \emph{content}: the field \code{ground} is now
refuted outright (\cref{thm:ropeLaw_universal_refuted}), not merely by the self-destruction of the
form. This is the same prohibition that forbids a perpetual engine on $\N$
(\code{no\_infinite\_descent}) and, read through Landauer, forbids extracting the cause from inside
--- knowledge is physical.

\begin{theorem}[verification, not derivation]\label{thm:collatz_verification_not_derivation}
\green{}\ A conjunction of three links: (1) infinite internal extraction of the rope law is impossible,
$\neg\,\mathrm{InternalKnowledgeOfCollatzCause}$; (2) the only internal trace of a counterexample is a
perpetual engine (\cref{thm:nonHalting_carries_perpetual_engine}); (3) hence the solution is either
unknowable from inside, or accessible only by \emph{verification} of an orbit.
\leantag{collatz\_verification\_not\_derivation}
\end{theorem}

The harnesses over $n\in[2,200000]$ with a divergence guard at $10^5$ steps are precisely the ``how
far'': no finite window can \emph{prove} the universal law --- it can only fail to find a refutation
within its horizon. And in the limit the universal law is false anyway, not because an orbit diverges
but because the law demanded more than is true; verifying convergence and verifying the law are not
the same thing.

\subsection{The decree taken, then withdrawn; open status}

The rope law was accepted \emph{by decree} --- as a fourth field
$\code{collatzBoundary}:\forall n\ge 1,\ \mathrm{RopeCountingLaw}(n)$ of the single axiom \axiom{},
its heuristic sign favourable and no forged refutation then known. Its price was disclosed at once:
only the arrow ``boundary $\Rightarrow$ conjecture'' was proven, so the decree might \emph{overpay}.
Two tripwires were planted alongside it --- intentional explosion detectors of the form ``a
refutation of the law $\Rightarrow \mathrm{False}$ exactly here''. When
\cref{thm:ropeLaw_universal_refuted} was forged, the tripwire fired: keeping the field would make the
axiom inconsistent. The field and its tripwires left the axiom together, and the boundaries are three
again (twins, Riemann, Navier--Stokes). The taint does not grow; after the removal it returned to
$47$. This closes the Collatz trilemma exactly as Yang--Mills was closed, by a forged refutation
rather than by a proof.

\begin{theorem}[open status of Collatz]\label{thm:collatz_open_status}
\green{}\ The epistemic status of Collatz after the decree's fall (the mirror of
\code{pnp\_locked\_behind\_engine\_status}, now \emph{without} the decree conjunct) is a conjunction of
four: $\neg\,\forall n\ge 1,\ \mathrm{RopeCountingLaw}(n)$ (the universal law is refuted) $\wedge$
$\neg\,\mathrm{InternalKnowledgeOfCollatzCause}$ (internal knowledge is impossible) $\wedge$
$\bigl(\forall n\ge 1,\ \mathrm{RopeCountingLaw}(n)\Rightarrow\exists K,\ \mathrm{iter}\,K\,n=1\bigr)$
(the per-$n$ law entails halting, conditional) $\wedge$
$\bigl(\forall n,\ (\forall K,\ \mathrm{iter}\,K\,n\neq 1)\Rightarrow \exists j,\
\mathrm{NonDescendingOrbit}(\mathrm{iter}\,j\,n)\wedge\forall K,\ \mathrm{iter}\,K\,(\mathrm{iter}\,j\,n)\neq 1\bigr)$
(a refutation would build a perpetual engine). \leantag{collatz\_open\_status}
\end{theorem}

This is not G\"odelian independence --- we do not claim Collatz is unprovable in Peano arithmetic.
What is proven, inside our system and rigorously, is that both internal paths (refutation and
self-grounding) run into one and the same forbidden object, and the decree third path is
machine-closed. The conjecture itself remains open.

\begin{conjecture}[Collatz]\label{conj:collatz}
\red{}\ For every $n\ge 1$ there is $K$ with $T^{K}(n)=1$; equivalently, every orbit has no perpetual tail.
Open, and untouched by the decree's fall: the refuted law was strictly stronger than the
conjecture, so its fall says only that Collatz cannot be closed by \emph{such} a decree.
\end{conjecture}

The classical statement is due to Collatz and surveyed by Lagarias~\cite{lagarias1985,lagarias2010}.
The lesson of the section is methodological: a decree that fell honestly --- its price disclosed, its
tripwire armed, its fall recorded as a theorem --- is worth more than a decree that stands unexamined.
The machine that verifies any conjecture is itself the object it cannot surpass.
